\documentclass[aspectratio=169]{beamer}
\usepackage{beamerthemeAMU}
\usepackage{amsmath}
\usepackage{booktabs}
\usepackage{array}
\usepackage{listings}
\usepackage{hyperref}
\usepackage[T1]{fontenc}
\usepackage[utf8]{inputenc}

\lstset{basicstyle=\ttfamily\small,keywordstyle=\color{AMU_dk2}\bfseries,commentstyle=\itshape\color{gray},showstringspaces=false,columns=flexible}

\title{Large Language Models in Data Science}
\subtitle{Week 6: Hackathon -- Applied LLM Projects}
\author{Sebastian Mueller}
\institute{Aix-Marseille Universit\'e}
\date{2025-2026}

\begin{document}

\begin{frame}[plain]
  \titlepage
\end{frame}

\begin{frame}{Session Overview}
  \begin{columns}[T,onlytextwidth]
    \begin{column}{0.55\linewidth}
      \textbf{Morning (Hackathon briefing)}
      \begin{enumerate}
        \item Recap of course concepts (Weeks 1--5)
        \item Hackathon rules and grading
        \item Project scope: what fits into one day
        \item Team formation and logistics
      \end{enumerate}
    \end{column}
    \begin{column}{0.4\linewidth}
      \textbf{Full-day Hackathon}
      \begin{itemize}
        \item Choose a problem and dataset
        \item Design an end-to-end LLM / RAG pipeline
        \item Implement, test, and iterate
        \item Prepare a short live demo and summary
      \end{itemize}
    \end{column}
  \end{columns}
\end{frame}

\section{Hackathon Overview}

\begin{frame}{Why a Hackathon?}
  \begin{itemize}
    \item Move from \textbf{toy examples} to an \textbf{end-to-end project} in a realistic setting.
    \item Combine techniques from the whole course: tokenization, embeddings, HF models, prompting, classification, RAG.
    \item Practice \textbf{teamwork} around LLMs and data: dividing roles, coordinating experiments, merging ideas.
    \item Experience trade-offs between ambition and scope under strict time constraints.
  \end{itemize}
\end{frame}

\begin{frame}{Evaluation \& Final Grade}
  \begin{itemize}
    \item The hackathon is the \textbf{final evaluation} of the course.
    \item Final score out of 20 is based on four criteria (5 points each), as in the README:
      \begin{itemize}
        \item \textbf{Use Case \& Applications (5 pts)}: relevance, originality, potential future usefulness.
        \item \textbf{Code Cleanliness (5 pts)}: readability, organization, modularity, reproducibility.
        \item \textbf{Lecture Integration (5 pts)}: use of course concepts (prompting, RAG, embeddings, classification, etc.).
        \item \textbf{Presentation (5 pts)}: clarity, conciseness, visuals, explanation of design choices.
      \end{itemize}
    \item Your hackathon score is your \textbf{official course grade}.
  \end{itemize}
\end{frame}

\begin{frame}{Teams \& Logistics}
  \begin{itemize}
    \item Work in groups of \textbf{1 to 5} students (recommended: 3--4).
    \item Each team produces:
      \begin{itemize}
        \item a small but coherent \textbf{project repository or notebook(s)},
        \item a short \textbf{live demo} of the system (or a recorded run, if needed),
        \item a \textbf{5--7 minute presentation} covering problem, approach, and lessons learned.
      \end{itemize}
    \item Use any combination of tools seen in the course:
      \begin{itemize}
        \item Hugging Face, Gemini API, scikit-learn, sentence-transformers, simple RAG stacks.
      \end{itemize}
    \item You are encouraged to reuse and adapt code from the labs (with proper structuring and documentation).
  \end{itemize}
\end{frame}

\section{Project Guidelines}

\begin{frame}{What Makes a Good Project?}
  \begin{itemize}
    \item \textbf{Clear problem statement}: for whom is the system built, and what decision does it help?
    \item \textbf{Grounded in data}: use real or realistic data relevant to AMU, Marseille, or data science practice.
    \item \textbf{Course concepts in action}:
      \begin{itemize}
        \item thoughtful prompts, embeddings, RAG, or classification models,
        \item not just a UI wrapped around an API.
      \end{itemize}
    \item \textbf{End-to-end story}: ingestion $\rightarrow$ processing / retrieval $\rightarrow$ model $\rightarrow$ output.
    \item \textbf{Simple but solid}: small scope, but robust, well-structured, and easy to demo.
  \end{itemize}
\end{frame}

\begin{frame}{Scope for One Day}
  \begin{itemize}
    \item Aim for a \textbf{minimal viable assistant}, not a full product.
    \item Prefer:
      \begin{itemize}
        \item one or two well-implemented workflows,
        \item over many incomplete features.
      \end{itemize}
    \item Reuse:
      \begin{itemize}
        \item lab pipelines (classification, RAG) as starting points,
        \item prompts and templates tested earlier in the course.
      \end{itemize}
    \item Keep infrastructure light:
      \begin{itemize}
        \item local notebooks, small datasets, API calls,
        \item no need for full web deployment if time is short.
      \end{itemize}
  \end{itemize}
\end{frame}

\section{Example Project Ideas}

\begin{frame}{Idea 1: AMU Data Science Programme Assistant}
  \begin{itemize}
    \item \textbf{Goal}: help prospective or current students explore the MAS / Data Science track at AMU.
    \item \textbf{Data}:
      \begin{itemize}
        \item official programme pages (like those used in the Week 5 lab),
        \item additional documents: FAQs, course descriptions, internship info.
      \end{itemize}
    \item \textbf{Pipeline}:
      \begin{itemize}
        \item ingest and chunk web pages / PDFs,
        \item build a BM25 + embedding index with metadata (campus, year, language),
        \item use a RAG chain to answer questions about courses, prerequisites, careers.
      \end{itemize}
    \item \textbf{Extensions}:
      \begin{itemize}
        \item multilingual Q\&A (French / English),
        \item citations and links back to source sections,
        \item a small evaluation set of typical student questions.
      \end{itemize}
  \end{itemize}
\end{frame}

\begin{frame}{Idea 2: Data Science EDA Copilot}
  \begin{itemize}
    \item \textbf{Goal}: assist a data scientist in exploring a tabular dataset (e.g., churn, housing, or a public Kaggle dataset).
    \item \textbf{Workflow}:
      \begin{itemize}
        \item user provides a CSV and a question (e.g., ``What drives churn in this dataset?''),
        \item the system generates Python code snippets for plots / tables using prompting techniques from Week 3,
        \item execute the code, collect results, and ask the LLM to summarise findings in natural language.
      \end{itemize}
    \item \textbf{Course concepts}:
      \begin{itemize}
        \item prompt templates for code generation and summarisation,
        \item structured outputs (e.g., JSON with ``steps'', ``code``, ``insights''),
        \item maybe a simple classification model from Week 4 as a baseline.
      \end{itemize}
    \item \textbf{Focus}:
      \begin{itemize}
        \item reproducibility (scripts/notebooks can be re-run),
        \item clear separation between data, code, and LLM outputs.
      \end{itemize}
  \end{itemize}
\end{frame}

\begin{frame}{Idea 3: Student Support Inbox Triage}
  \begin{itemize}
    \item \textbf{Goal}: help route incoming student emails to the right service and provide draft replies.
    \item \textbf{Data}:
      \begin{itemize}
        \item a synthetic dataset of student enquiries (motivated by Week 4 intent classification),
        \item labels such as \texttt{timetable}, \texttt{exam}, \texttt{administrative}, \texttt{internship}, \texttt{other}.
      \end{itemize}
    \item \textbf{Pipeline}:
      \begin{itemize}
        \item baseline classifier (embeddings + logistic regression or zero-shot classification),
        \item optional RAG component with a small FAQ corpus,
        \item LLM prompting to generate a draft response and a recommended action.
      \end{itemize}
    \item \textbf{Evaluation}:
      \begin{itemize}
        \item accuracy of routing to labels (using a held-out test set),
        \item qualitative assessment of suggested replies (are they precise and polite?).
      \end{itemize}
  \end{itemize}
\end{frame}

\section{Planning Your Project}

\begin{frame}{Checklist for Today}
  \begin{itemize}
    \item \textbf{1. Define the problem}:
      \begin{itemize}
        \item target user, key questions, success criteria.
      \end{itemize}
    \item \textbf{2. Select data \& models}:
      \begin{itemize}
        \item datasets, external documents, HF models, LLM endpoints.
      \end{itemize}
    \item \textbf{3. Design the pipeline}:
      \begin{itemize}
        \item retrieval / preprocessing, model calls, post-processing.
      \end{itemize}
    \item \textbf{4. Implement the minimal version}:
      \begin{itemize}
        \item one or two key workflows working end-to-end.
      \end{itemize}
    \item \textbf{5. Polish \& present}:
      \begin{itemize}
        \item clean up code, add README, prepare a concise demo.
      \end{itemize}
  \end{itemize}
\end{frame}

\begin{frame}{Final Remarks}
  \begin{itemize}
    \item Use this day to \textbf{experiment} and \textbf{connect ideas} from the course.
    \item Keep the scope realistic, but do not hesitate to be creative in your design.
    \item Remember the evaluation criteria: clarity of use case, clean code, integration of course content, and a clear presentation.
    \item Most importantly: have fun building with LLMs and data.
  \end{itemize}
\end{frame}

\end{document}

